%!TeX root=../emmatop.tex
\chapter[Chapter \thechapter]{}
\lettrine[lraise=0.3]{T}{he} weather continued much the same all the following morning; and the same loneliness, and the same melancholy, seemed to reign at Hartfield—but in the afternoon it cleared; the wind changed into a softer quarter; the clouds were carried off; the sun appeared; it was summer again. With all the eagerness which such a transition gives, Emma resolved to be out of doors as soon as possible. Never had the exquisite sight, smell, sensation of nature, tranquil, warm, and brilliant after a storm, been more attractive to her. She longed for the serenity they might gradually introduce; and on Mr Perry's coming in soon after dinner, with a disengaged hour to give her father, she lost no time in hurrying into the shrubbery.—There, with spirits freshened, and thoughts a little relieved, she had taken a few turns, when she saw Mr Knightley passing through the garden door, and coming towards her.—It was the first intimation of his being returned from London. She had been thinking of him the moment before, as unquestionably sixteen miles distant.—There was time only for the quickest arrangement of mind. She must be collected and calm. In half a minute they were together. The <How d'ye do's> were quiet and constrained on each side. She asked after their mutual friends; they were all well.—When had he left them?—Only that morning. He must have had a wet ride.—Yes.—He meant to walk with her, she found. <He had just looked into the dining-room, and as he was not wanted there, preferred being out of doors.>—She thought he neither looked nor spoke cheerfully; and the first possible cause for it, suggested by her fears, was, that he had perhaps been communicating his plans to his brother, and was pained by the manner in which they had been received.

They walked together. He was silent. She thought he was often looking at her, and trying for a fuller view of her face than it suited her to give. And this belief produced another dread. Perhaps he wanted to speak to her, of his attachment to Harriet; he might be watching for encouragement to begin.—She did not, could not, feel equal to lead the way to any such subject. He must do it all himself. Yet she could not bear this silence. With him it was most unnatural. She considered—resolved—and, trying to smile, began—

<You have some news to hear, now you are come back, that will rather surprize you.>

<Have I\@?> said he quietly, and looking at her; <of what nature?>

<Oh! the best nature in the world—a wedding.>

After waiting a moment, as if to be sure she intended to say no more, he replied,

<If you mean Miss Fairfax and Frank Churchill, I have heard that already.>

<How is it possible?> cried Emma, turning her glowing cheeks towards him; for, while she spoke, it occurred to her that he might have called at Mrs Goddard's in his way.

<I had a few lines on parish business from Mr Weston this morning, and at the end of them he gave me a brief account of what had happened.>

Emma was quite relieved, and could presently say, with a little more composure,

<You probably have been less surprized than any of us, for you have had your suspicions.—I have not forgotten that you once tried to give me a caution.—I wish I had attended to it—but—(with a sinking voice and a heavy sigh) I seem to have been doomed to blindness.>

For a moment or two nothing was said, and she was unsuspicious of having excited any particular interest, till she found her arm drawn within his, and pressed against his heart, and heard him thus saying, in a tone of great sensibility, speaking low,

<Time, my dearest Emma, time will heal the wound.—Your own excellent sense—your exertions for your father's sake—I know you will not allow yourself—.> Her arm was pressed again, as he added, in a more broken and subdued accent, <The feelings of the warmest friendship—Indignation—Abominable scoundrel!>—And in a louder, steadier tone, he concluded with, <He will soon be gone. They will soon be in Yorkshire. I am sorry for her. She deserves a better fate.>

Emma understood him; and as soon as she could recover from the flutter of pleasure, excited by such tender consideration, replied,

<You are very kind—but you are mistaken—and I must set you right.— I am not in want of that sort of compassion. My blindness to what was going on, led me to act by them in a way that I must always be ashamed of, and I was very foolishly tempted to say and do many things which may well lay me open to unpleasant conjectures, but I have no other reason to regret that I was not in the secret earlier.>

<Emma!> cried he, looking eagerly at her, <are you, indeed?>—but checking himself—<No, no, I understand you—forgive me—I am pleased that you can say even so much.—He is no object of regret, indeed! and it will not be very long, I hope, before that becomes the acknowledgment of more than your reason.—Fortunate that your affections were not farther entangled!—I could never, I confess, from your manners, assure myself as to the degree of what you felt—I could only be certain that there was a preference—and a preference which I never believed him to deserve.—He is a disgrace to the name of man.—And is he to be rewarded with that sweet young woman?—Jane, Jane, you will be a miserable creature.>

<Mr Knightley,> said Emma, trying to be lively, but really confused—<I am in a very extraordinary situation. I cannot let you continue in your error; and yet, perhaps, since my manners gave such an impression, I have as much reason to be ashamed of confessing that I never have been at all attached to the person we are speaking of, as it might be natural for a woman to feel in confessing exactly the reverse.—But I never have.>

He listened in perfect silence. She wished him to speak, but he would not. She supposed she must say more before she were entitled to his clemency; but it was a hard case to be obliged still to lower herself in his opinion. She went on, however.

<I have very little to say for my own conduct.—I was tempted by his attentions, and allowed myself to appear pleased.—An old story, probably—a common case—and no more than has happened to hundreds of my sex before; and yet it may not be the more excusable in one who sets up as I do for Understanding. Many circumstances assisted the temptation. He was the son of Mr Weston—he was continually here—I always found him very pleasant—and, in short, for (with a sigh) let me swell out the causes ever so ingeniously, they all centre in this at last—my vanity was flattered, and I allowed his attentions. Latterly, however—for some time, indeed—I have had no idea of their meaning any thing.—I thought them a habit, a trick, nothing that called for seriousness on my side. He has imposed on me, but he has not injured me. I have never been attached to him. And now I can tolerably comprehend his behaviour. He never wished to attach me. It was merely a blind to conceal his real situation with another.—It was his object to blind all about him; and no one, I am sure, could be more effectually blinded than myself—except that I was not blinded—that it was my good fortune—that, in short, I was somehow or other safe from him.>

She had hoped for an answer here—for a few words to say that her conduct was at least intelligible; but he was silent; and, as far as she could judge, deep in thought. At last, and tolerably in his usual tone, he said,

<I have never had a high opinion of Frank Churchill.—I can suppose, however, that I may have underrated him. My acquaintance with him has been but trifling.—And even if I have not underrated him hitherto, he may yet turn out well.—With such a woman he has a chance.—I have no motive for wishing him ill—and for her sake, whose happiness will be involved in his good character and conduct, I shall certainly wish him well.>

<I have no doubt of their being happy together,> said Emma; <I believe them to be very mutually and very sincerely attached.>

<He is a most fortunate man!> returned Mr Knightley, with energy. <So early in life—at three-and-twenty—a period when, if a man chuses a wife, he generally chuses ill. At three-and-twenty to have drawn such a prize! What years of felicity that man, in all human calculation, has before him!—Assured of the love of such a woman—the disinterested love, for Jane Fairfax's character vouches for her disinterestedness; every thing in his favour,—equality of situation—I mean, as far as regards society, and all the habits and manners that are important; equality in every point but one—and that one, since the purity of her heart is not to be doubted, such as must increase his felicity, for it will be his to bestow the only advantages she wants.—A man would always wish to give a woman a better home than the one he takes her from; and he who can do it, where there is no doubt of her regard, must, I think, be the happiest of mortals.—Frank Churchill is, indeed, the favourite of fortune. Every thing turns out for his good.—He meets with a young woman at a watering-place, gains her affection, cannot even weary her by negligent treatment—and had he and all his family sought round the world for a perfect wife for him, they could not have found her superior.—His aunt is in the way.—His aunt dies.—He has only to speak.—His friends are eager to promote his happiness.—He had used every body ill—and they are all delighted to forgive him.—He is a fortunate man indeed!>

<You speak as if you envied him.>

<And I do envy him, Emma. In one respect he is the object of my envy.>

Emma could say no more. They seemed to be within half a sentence of Harriet, and her immediate feeling was to avert the subject, if possible. She made her plan; she would speak of something totally different—the children in Brunswick Square; and she only waited for breath to begin, when Mr Knightley startled her, by saying,

<You will not ask me what is the point of envy.—You are determined, I see, to have no curiosity.—You are wise—but I cannot be wise. Emma, I must tell you what you will not ask, though I may wish it unsaid the next moment.>

<Oh! then, don't speak it, don't speak it,> she eagerly cried. <Take a little time, consider, do not commit yourself.>

<Thank you,> said he, in an accent of deep mortification, and not another syllable followed.

Emma could not bear to give him pain. He was wishing to confide in her—perhaps to consult her;—cost her what it would, she would listen. She might assist his resolution, or reconcile him to it; she might give just praise to Harriet, or, by representing to him his own independence, relieve him from that state of indecision, which must be more intolerable than any alternative to such a mind as his.—They had reached the house.

<You are going in, I suppose?> said he.

<No,>—replied Emma—quite confirmed by the depressed manner in which he still spoke—<I should like to take another turn. Mr Perry is not gone.> And, after proceeding a few steps, she added—<I stopped you ungraciously, just now, Mr Knightley, and, I am afraid, gave you pain.—But if you have any wish to speak openly to me as a friend, or to ask my opinion of any thing that you may have in contemplation—as a friend, indeed, you may command me.—I will hear whatever you like. I will tell you exactly what I think.>

<As a friend!>—repeated Mr Knightley.—<Emma, that I fear is a word—No, I have no wish—Stay, yes, why should I hesitate?—I have gone too far already for concealment.—Emma, I accept your offer—Extraordinary as it may seem, I accept it, and refer myself to you as a friend.—Tell me, then, have I no chance of ever succeeding?>

\begin{figure}[tbph]
\centering
\includegraphics[width=.9\linewidth]{49question}
\caption{He stopped to look the question}
\end{figure}

He stopped in his earnestness to look the question, and the expression of his eyes overpowered her.

<My dearest Emma,> said he, <for dearest you will always be, whatever the event of this hour's conversation, my dearest, most beloved Emma—tell me at once. Say <No,> if it is to be said.>—She could really say nothing.—<You are silent,> he cried, with great animation; <absolutely silent! at present I ask no more.>

Emma was almost ready to sink under the agitation of this moment. The dread of being awakened from the happiest dream, was perhaps the most prominent feeling.

<I cannot make speeches, Emma:> he soon resumed; and in a tone of such sincere, decided, intelligible tenderness as was tolerably convincing.—<If I loved you less, I might be able to talk about it more. But you know what I am.—You hear nothing but truth from me.—I have blamed you, and lectured you, and you have borne it as no other woman in England would have borne it.—Bear with the truths I would tell you now, dearest Emma, as well as you have borne with them. The manner, perhaps, may have as little to recommend them. God knows, I have been a very indifferent lover.—But you understand me.—Yes, you see, you understand my feelings—and will return them if you can. At present, I ask only to hear, once to hear your voice.>

While he spoke, Emma's mind was most busy, and, with all the wonderful velocity of thought, had been able—and yet without losing a word—to catch and comprehend the exact truth of the whole; to see that Harriet's hopes had been entirely groundless, a mistake, a delusion, as complete a delusion as any of her own—that Harriet was nothing; that she was every thing herself; that what she had been saying relative to Harriet had been all taken as the language of her own feelings; and that her agitation, her doubts, her reluctance, her discouragement, had been all received as discouragement from herself.—And not only was there time for these convictions, with all their glow of attendant happiness; there was time also to rejoice that Harriet's secret had not escaped her, and to resolve that it need not, and should not.—It was all the service she could now render her poor friend; for as to any of that heroism of sentiment which might have prompted her to entreat him to transfer his affection from herself to Harriet, as infinitely the most worthy of the two—or even the more simple sublimity of resolving to refuse him at once and for ever, without vouchsafing any motive, because he could not marry them both, Emma had it not. She felt for Harriet, with pain and with contrition; but no flight of generosity run mad, opposing all that could be probable or reasonable, entered her brain. She had led her friend astray, and it would be a reproach to her for ever; but her judgment was as strong as her feelings, and as strong as it had ever been before, in reprobating any such alliance for him, as most unequal and degrading. Her way was clear, though not quite smooth.—She spoke then, on being so entreated.—What did she say?—Just what she ought, of course. A lady always does.—She said enough to shew there need not be despair—and to invite him to say more himself. He had despaired at one period; he had received such an injunction to caution and silence, as for the time crushed every hope;—she had begun by refusing to hear him.—The change had perhaps been somewhat sudden;—her proposal of taking another turn, her renewing the conversation which she had just put an end to, might be a little extraordinary!—She felt its inconsistency; but Mr Knightley was so obliging as to put up with it, and seek no farther explanation.

Seldom, very seldom, does complete truth belong to any human disclosure; seldom can it happen that something is not a little disguised, or a little mistaken; but where, as in this case, though the conduct is mistaken, the feelings are not, it may not be very material.—Mr Knightley could not impute to Emma a more relenting heart than she possessed, or a heart more disposed to accept of his.

He had, in fact, been wholly unsuspicious of his own influence. He had followed her into the shrubbery with no idea of trying it. He had come, in his anxiety to see how she bore Frank Churchill's engagement, with no selfish view, no view at all, but of endeavouring, if she allowed him an opening, to soothe or to counsel her.—The rest had been the work of the moment, the immediate effect of what he heard, on his feelings. The delightful assurance of her total indifference towards Frank Churchill, of her having a heart completely disengaged from him, had given birth to the hope, that, in time, he might gain her affection himself;—but it had been no present hope—he had only, in the momentary conquest of eagerness over judgment, aspired to be told that she did not forbid his attempt to attach her.—The superior hopes which gradually opened were so much the more enchanting.—The affection, which he had been asking to be allowed to create, if he could, was already his!—Within half an hour, he had passed from a thoroughly distressed state of mind, to something so like perfect happiness, that it could bear no other name.

Her change was equal.—This one half-hour had given to each the same precious certainty of being beloved, had cleared from each the same degree of ignorance, jealousy, or distrust.—On his side, there had been a long-standing jealousy, old as the arrival, or even the expectation, of Frank Churchill.—He had been in love with Emma, and jealous of Frank Churchill, from about the same period, one sentiment having probably enlightened him as to the other. It was his jealousy of Frank Churchill that had taken him from the country.—The Box Hill party had decided him on going away. He would save himself from witnessing again such permitted, encouraged attentions.—He had gone to learn to be indifferent.—But he had gone to a wrong place. There was too much domestic happiness in his brother's house; woman wore too amiable a form in it; Isabella was too much like Emma—differing only in those striking inferiorities, which always brought the other in brilliancy before him, for much to have been done, even had his time been longer.—He had stayed on, however, vigorously, day after day—till this very morning's post had conveyed the history of Jane Fairfax.—Then, with the gladness which must be felt, nay, which he did not scruple to feel, having never believed Frank Churchill to be at all deserving Emma, was there so much fond solicitude, so much keen anxiety for her, that he could stay no longer. He had ridden home through the rain; and had walked up directly after dinner, to see how this sweetest and best of all creatures, faultless in spite of all her faults, bore the discovery.

He had found her agitated and low.—Frank Churchill was a villain.— He heard her declare that she had never loved him. Frank Churchill's character was not desperate.—She was his own Emma, by hand and word, when they returned into the house; and if he could have thought of Frank Churchill then, he might have deemed him a very good sort of fellow.